\subsection{对称算子的自伴扩张}

在本节中，我们打算对复希尔伯特空间上对称算子的自伴扩张进行分类，特别是找出保证这类算子的自伴扩张存在的条件。

设 E 是复希尔伯特空间。对 E 上每个部分算子 u，本节中我们将把 \(u + i\) 与 \(u - i\) 分别写作 \(u + i1_{E}\) 与 \(u - i1_{E}\) 的简写。

设 u 是 E 上稠定的闭算子。我们回顾（参见 IV，第 230 页定义 4 及 IV，第 230 页命题 4）：我们用 \(E_{u}\) 表示赋有埃尔米特形式的希尔伯特空间 dom(u)

\[
(x \mid y) _ {u} = \langle x \mid y \rangle + \langle u (x) \mid u (y) \rangle .
\]

我们用 \(\|x\|_{u}\) 表示希尔伯特空间 \(E_{u}\) 中元素 x 的范数。dom(u) 的每个在 E 中闭的子空间都在 \(E_{u}\) 中闭（IV，第 230 页，注 1）。

定义 10. --- 设 u 是复希尔伯特空间 E 上的闭对称算子。设 \(\mathrm{E}_{+} = \mathrm{Ker}(u^{*} - i)\)，\(\mathrm{E}_{-} = \mathrm{Ker}(u^{*} + i)\)。子空间偶 \((\mathrm{E}_{+}, \mathrm{E}_{-})\)（\(\mathrm{dom}(u^{*})\) 的）称为 u 的亏偶。子空间 \(E_{+}\) 与 \(E_{-}\) 是 E 的闭子空间（IV，第 229 页，命题 3）。这些子空间的希尔伯特维数所成的偶 \((\dim(\mathrm{E}_{+}), \dim(\mathrm{E}_{-}))\) 称为 u 的亏指标。

命题 25. --- 设 u 是复希尔伯特空间 E 上的闭对称算子。我们有

\[
\begin{array}{l} \operatorname{Ker} (u ^ {*} - i) ^ {\circ} = \operatorname{Im} (u + i), \\ \operatorname{Ker} (u ^ {*} + i) ^ {\circ} = \operatorname{Im} (u - i). \end{array}
\]

由 IV，第 236 页命题 7，c)，只需证明 \(u + i\)（相应地，\(u - i\)）的像在 E 中闭。由于 \(u\) 对称，我们有 \(\langle (u + i)(x) \mid (u + i)(y) \rangle = (x \mid y)_u\)，对一切 \(x\) 与 \(y\)（\(\text{dom}(u)\) 中的）。于是映射 \(x \mapsto u(x) + ix\)（从 \(E_u\) 到 E 中的）是等距的。它的像在 E 中闭（I，第 107 页，引理 8）。由于这也是 \(u + i\) 的像，故后者在 E 中闭。同理，\(u - i\) 的像在 E 中闭。

引理 7. --- 设 u 是复希尔伯特空间 E 上的闭对称算子，v 是 u 的闭对称扩张。u 的定义域是希尔伯特空间 \(E_{v}\) 的闭子空间。

由于 v 是 u 的扩张，我们有 \((x \mid y)_{v} = (x \mid y)_{u}\)，对 \(\text{dom}(u)\) 中的 x 与 y。于是 \(E_{u}\) 到 \(E_{v}\) 中的典范单射是等距的，结论由此得出（I，第 107 页，引理 8）。

命题 26. --- 设 u 是复希尔伯特空间 E 上的闭对称算子。设 \((\mathrm{E}_{+}, \mathrm{E}_{-})\) 是 u 的亏偶。空间 \(E_{+}\)、\(E_{-}\) 与 dom(u) 是 \(E_{u^{*}}\) 中两两正交的闭子空间，其希尔伯特直和等于 \(E_{u^{*}}\)。

由于 u 对称，我们有 \(u \subset u^{*}\)，故空间 \(\text{dom}(u)\) 在 \(E_{u^{*}}\) 中闭（引理 7）。子空间 \(E_{+}\) 与 \(E_{-}\) 在 E 中闭且包含于 \(\text{dom}(u^{*})\)，故在 \(E_{u^{*}}\) 中闭。

设 \(x \in \mathrm{E}_{+}\)。我们有 \(u^{*}(x) = ix\)，故 \(\langle x | u(y) \rangle = -i\langle x | y \rangle\) 对一切 \(y \in \operatorname{dom}(u)\) 成立。于是，对每个 \(y \in \operatorname{dom}(u)\)，有

\[
(x \mid y) _ {u ^ {*}} = \langle x \mid y \rangle + \langle u ^ {*} (x) \mid u ^ {*} (y) \rangle = i (\langle x \mid u (y) \rangle - \langle x \mid u ^ {*} (y) \rangle) = 0
\]

因为 u 对称。因此空间 \(E_{+}\) 与 \(\operatorname{dom}(u)\) 在 \(E_{u^{*}}\) 中正交。同理，\(E_{-}\) 与 \(\operatorname{dom}(u)\) 在 \(E_{u^{*}}\) 中正交。

设 \(x \in \mathrm{E}_{+}\)，\(y \in \mathrm{E}_{-}\)。则

\[
(x \mid y) _ {u ^ {*}} = \langle x \mid y \rangle + \langle u ^ {*} (x) \mid u ^ {*} (y) \rangle = \langle x \mid y \rangle - \langle x \mid y \rangle = 0,
\]

故 \(\mathrm{E}_{+}\) 与 \(\mathrm{E}_{-}\) 在 \(\mathrm{E}_{u^{*}}\) 中正交。

设 \(x \in \mathrm{E}_{u^*}\) 与闭子空间 \(\operatorname{dom}(u) \oplus \mathrm{E}_+ \oplus \mathrm{E}_-\) 正交。对每个 \(y \in \operatorname{dom}(u)\)，有 \(u^*(y) = u(y)\)（因为 \(u \subset u^*\)），于是

\[
0 = (x \mid y) _ {u ^ {*}} = \langle x \mid y \rangle + \langle u ^ {*} (x) \mid u ^ {*} (y) \rangle = \langle x \mid y \rangle + \langle u ^ {*} (x) \mid u (y) \rangle .
\]

这蕴含 \(z = u^{*}(x)\) 属于 \(u^{*}\) 的定义域且满足 \(u^{*}(z) = -x\)。令 \(x_{-} = z - ix\)。它是 \(\text{dom}(u^{*})\) 中满足 \(u^{*}(x_{-}) = -x - iz = -ix_{-}\) 的元素。故 \(x_{-}\) 属于 \(E_{-}\)。但是，对每个 \(w \in E_{-}\)，有

\[
\begin{array}{c} - i \langle x _ {-} | w \rangle = - i \langle z | w \rangle + \langle x | w \rangle \\ = \langle u ^ {*} (x) | u ^ {*} (w) \rangle + \langle x | w \rangle = (x | w) _ {u ^ {*}} = 0. \end{array}
\]

取 \(w = x_{-}\)，便得 \(x_{-} = 0\)，即 \(z = ix\)；由于 \(z = u^{*}(x)\)，我们有 \(x \in \mathrm{E}_{+}\)，从而 \(x = 0\)，因为 \(x\) 与 \(\mathrm{E}_{+}\) 在 \(\mathrm{E}_{u^{*}}\) 中正交。这就完成了直和 \(\operatorname{dom}(u) \oplus \mathrm{E}_{+} \oplus \mathrm{E}_{-}\) 等于 \(\mathrm{E}_{u^{*}}\) 的证明。

推论。--- 设 u 是 E 上的闭对称算子。则 u 自伴，当且仅当 u 的亏指标等于 (0,0)。

事实上，对称算子 \(u\) 自伴当且仅当 \(\text{dom}(u^*) = \text{dom}(u)\)；而命题 26 表明 \(\text{dom}(u^*)\) 是 \(\text{dom}(u)\) 与 \(u\) 的亏子空间的希尔伯特直和。

\pandocbounded{\includegraphics[keepaspectratio]{images/part002_f39b57bf.jpg}}

设 \(u\) 是闭对称部分算子。\(u\) 的谱包含于 \(\mathbf{R}\)，当且仅当 \(u\) 自伴。事实上，我们知道 \(\mathrm{Sp}(u) \subset \mathbf{R}\) 当 \(u\) 自伴时成立（IV，第 248 页，命题 17）；若 \(u\) 不自伴，则子空间 \(\mathrm{Ker}(u^{*} + i)\) 与 \(\mathrm{Ker}(u^{*} - i)\) 之一非零，于是 \(u + i\) 或 \(u - i\) 的像是 E 的真子空间（IV，第 256 页，命题 25），从而 \(i\) 或 \(-i\) 属于 \(\mathrm{Sp}(u)\)。

定义 11. --- 设 u 是复希尔伯特空间 E 上的闭对称算子。u 的边界条件是从 \(\operatorname{Ker}(u^{*} - i)\) 到 \(\operatorname{Ker}(u^{*} + i)\) 中的一个部分等距线性映射（EVT，V，第 41 页，定义 3）。

设 \(u\) 是复希尔伯特空间 E 上的闭对称算子，\(b\) 是 \(u\) 的一个边界条件。设 \(I = \text{Ker}(b)^{\circ}\) 为 \(b\) 的初始子空间。我们用 \(u_b\) 表示 \(u^*\) 到 \(\text{dom}(u^*)\) 的如下子空间上的缩减：该子空间是 \(\text{dom}(u)\) 与一个图像的直和，该图像位于 \(E_+ \oplus E_-\) 中，是 \(b\) 在 I 上的限制的图像。由于 \(E_{u^*} = \text{dom}(u) \oplus E_+ \oplus E_-\)（命题 26），部分算子 \(u_b\) 是完全确定的。

换句话说，\(u_{b}\) 的定义域是这样的元素 \(x \in E\) 所成的空间：它们形如 \(x = x_{0} + y + b(y)\)，其中 \(x_{0} \in \text{dom}(u)\)，y 属于 b 的初始子空间。此时我们有 \(u_{b}(x_{0} + y + b(y)) = u(x_{0}) + iy - ib(y)\)。

定理 1（冯·诺伊曼）。--- 设 u 是复希尔伯特空间 E 上的闭对称算子，\((\mathrm{E}_{+}, \mathrm{E}_{-})\) 是它的亏偶。

\begin{enumerate}
\def\labelenumi{\alph{enumi})}
\item
  对 u 的每个边界条件 b，部分算子 \(u_{b}\) 是 u 的闭对称扩张；
\item
  映射 \(b \mapsto u_{b}\) 是从 u 的边界条件集合到 u 的闭对称扩张集合上的双射；
\item
  对 u 的每个边界条件 b，我们有
\end{enumerate}

\[
\begin{array}{l} \operatorname{Im} (u _ {b} + i) = \operatorname{Im} (u + i) \oplus \operatorname{Im} (b) \\ \operatorname{Im} (u _ {b} - i) = \operatorname{Im} (u - i) \oplus \operatorname{Ker} (b) ^ {\circ}; \end{array}
\]

\begin{enumerate}
\def\labelenumi{\alph{enumi})}
\setcounter{enumi}{3}
\tightlist
\item
  对 u 的每个边界条件 b，\(u_{b}\) 的亏偶是 \((\operatorname{Ker}(b), \operatorname{Ker}(b^{*}))\)。
\end{enumerate}

设 \(b\) 是 \(u\) 的一个边界条件，\(\mathrm{I} \subset \mathrm{E}_{+}\) 是它的初始子空间。\(u^{*}\) 到 \(\Gamma_{b|I} \subset \mathrm{E}_{+} \oplus \mathrm{E}_{-}\) 的限制是由 \(x + b(x) \mapsto ix - ib(x)\)（\(x \in I\)）定义的连续线性映射。\(u_{b}\) 的图像是 u 的图像与该线性映射的图像的直和；因此它是闭的。闭算子 \(u_{b}\) 是 u 的扩张且满足 \(u_{b} \subset u^{*}\)。

设 \(\gamma_{1}\) 与 \(\eta_{1}\) 是 I 的元素。考虑元素

\[
\gamma = \gamma_ {1} + b (\gamma_ {1}), \quad \eta = \eta_ {1} + b (\eta_ {1})
\]

它们是图像 \(\Gamma_{b|I}\) 中的元素。我们计算

\[
\begin{array}{c} \langle \gamma | u ^ {*} (\eta) \rangle = i \big (\langle b (\gamma_ {1}) | \eta_ {1} \rangle - \langle \gamma_ {1} | b (\eta_ {1}) \rangle \big) + i \big (\langle \gamma_ {1} | \eta_ {1} \rangle - \langle b (\gamma_ {1}) | b (\eta_ {1}) \rangle \big) \\ = i \big (\langle b (\gamma_ {1}) | \eta_ {1} \rangle - \langle \gamma_ {1} | b (\eta_ {1}) \rangle \big) \end{array}
\]

因为 b 在 I 上等距。由此我们推出

\[
\langle u ^ {*} (\gamma) \mid \eta \rangle = \langle \gamma \mid u ^ {*} (\eta) \rangle\tag{12}
\]

对一切 \(\gamma\) 与 \(\eta\)（在 \(\Gamma_{b|I}\) 中）成立。

设 x 与 y 是 \(\operatorname{dom}(u)\) 的元素，\(\gamma\) 与 \(\eta\) 是 \(\Gamma_{b|I}\) 的元素。我们得到

\[
\begin{array}{c} \langle x + \gamma | u _ {b} (y + \eta) \rangle = \langle x | u (y) \rangle + \langle x | u ^ {*} (\eta) \rangle + \langle \gamma | u (y) \rangle + \langle \gamma | u ^ {*} (\eta) \rangle \\ = \langle u (x) | y \rangle + \langle u (x) | \eta \rangle + \langle u ^ {*} (\gamma) | y \rangle + \langle u ^ {*} (\gamma) | \eta \rangle \\ = \langle u _ {b} (x + \gamma) | y + \eta \rangle , \end{array}
\]

其中用到了 u 对称这一事实、\(\gamma\) 与 \(\eta\) 属于 \(u^{*}\) 的定义域这一事实，以及公式 (12)。于是算子 \(u_{b}\) 是 u 的闭对称扩张。断言 a) 得证。

映射 \(b \mapsto u_b\) 是单射。事实上，从 \(E_+\) 到 \(E_-\) 中的部分等距映射由其初始子空间及其在该子空间上的限制唯一确定。而 \(u_b\) 的定义域确定了 \(b\) 的初始子空间 I 以及 \(b\) 在 I 上限制的图像。

我们来证明映射 \(b \mapsto u_b\) 是满射。设 \(w\) 是 \(u\) 的闭对称扩张。我们有 \(w \subset w^* \subset u^*\)。于是 \(w\) 的定义域是 \(\text{dom}(u^*)\) 中包含 \(\text{dom}(u)\) 的子空间，且 \(w\) 是 \(u^*\) 到该子空间上的缩减。设 G 是 \(w\) 的定义域与 \(E_+ \oplus E_-\) 的交。它是 \(E_{u^*}\) 的闭子空间（引理 7），且由命题 26 有 \(\text{dom}(w) = \text{dom}(u) \oplus G\)。

设 \(x \in \mathrm{E}_{+}\)，\(y \in \mathrm{E}_{-}\) 满足 \(x + y \in \mathrm{G}\)。由于 \(w\) 对称且是 \(u^{*}\) 到 \(\operatorname{dom}(w)\) 上的缩减，我们有

\[
\langle x + y \mid u ^ {*} (x + y) \rangle = \langle u ^ {*} (x + y) \mid x + y \rangle .
\]

此等式等价于

\[
\begin{array}{c} \langle x | i x \rangle + \langle x | - i y \rangle + \langle y | i x \rangle + \langle y | - i y \rangle \\ = \langle i x | x \rangle + \langle i x | y \rangle + \langle - i y | x \rangle + \langle - i y | y \rangle , \end{array}
\]

亦即等价于 \(\|x\|^{2} = \|y\|^{2}\)。特别地，典范投影 \(p_{+}: G \to E_{+}\) 是单射，且若用 I 表示它的像，则空间 G 是等距映射 \(b_{0}\)（从 I 到 \(E_{-}\) 中的）的图像。对 I 中的 x 与 \(E_{-}\) 中满足 \(x + y \in G\) 的 y，我们有

\[
\begin{array}{r} \| x + y \| _ {u ^ {*}} ^ {2} = \| x \| _ {u ^ {*}} ^ {2} + \| y \| _ {u ^ {*}} ^ {2} = 2 (\| x \| ^ {2} + \| y \| ^ {2}) \\ = 4 \| x \| ^ {2} = 4 \| p _ {+} (x + y) \| ^ {2}. \end{array}
\]

于是，把 G 看作希尔伯特空间 \(E_{u^{*}}\) 的闭子空间时，映射 \(x \mapsto \frac{1}{2}p_{+}(x)\) 是从 G 到 I 上的等距映射。特别地，I 在 \(E_{+}\) 中闭（I，第 107 页，引理 8）。于是存在 u 的唯一的边界条件 b，以 I 为初始空间且在 I 上与 \(b_{0}\) 相合。我们有 \(\text{dom}(w) = \text{dom}(u_{b})\)，从而 \(w = u_{b}\)。这就建立了断言 b)。

我们来证明断言 \(c\)）与 \(d\)）。设 \(b\) 是 \(u\) 的一个边界条件，I 是它的初始子空间。对 \(x_0 \in \mathrm{dom}(u)\) 与 \(y \in \mathrm{I}\)，有

\[
u _ {b} (x _ {0} + y + b (y)) - i (x _ {0} + y + b (y)) = (u - i) (x _ {0}) - 2 i b (y).
\]

我们有 \(\operatorname{Im}(u - i) = \operatorname{E}_{-}^{\circ}\)（命题 25）。由于 \(b\) 的像包含在 \(\operatorname{E}_{-}\) 中，该公式表明 \(\operatorname{Im}(u_b - i) = \operatorname{Im}(u - i) \oplus \operatorname{Im}(b)\)。由同上所引，我们因此有 \(\operatorname{Ker}(u_b^* + i) = \operatorname{E}_{-} \cap \operatorname{Im}(b)^{\circ} = \operatorname{Ker}(b^*)\)。

类似地，可以验证 \(\operatorname{Im}(u_b + i) = \operatorname{Im}(u + i) \oplus \operatorname{I}\)，从而（同上所引）有 \(\operatorname{Ker}(u_b^* - i) = \operatorname{I}^\circ \cap \operatorname{E}_+ = \operatorname{Ker}(b)\)。

推论 1. --- 设 u 是复希尔伯特空间 E 上的闭对称算子。映射 \(b \mapsto u_b\) 诱导出从 \(\text{Ker}(u^* - i)\) 到 \(\text{Ker}(u^* + i)\) 上的等距同构集合与 u 的自伴扩张集合之间的双射。特别地，存在 u 的自伴扩张，当且仅当 u 的亏指标的两个分量相等。

这由命题 26 的推论与该定理以及 EVT，V，第 25 页，推论 2 得出。

推论 2. --- 设 E 是实希尔伯特空间。设 u 是 E 上的闭对称算子。部分算子 \(u_{(\mathbf{C})}\)（\(\mathrm{E}_{(\mathbf{C})}\) 上的）对称且闭，并且至少有一个自伴扩张。

部分算子 \(u_{(\mathbf{C})}\) 闭且对称（IV，第 228 页命题 2 与 IV，第 239 页引理 3）。同构 \(j\)（从 \(\mathrm{E}_{(\mathbf{C})}\) 到 \(\mathrm{E}_{(\mathbf{C})}\) 上的 R-线性同构）满足 \(j(z \otimes x) = \overline{z} \otimes x\)，对一切 \(z \in \mathbf{C}\) 与每个 \(x \in \mathrm{E}\)；它诱导出从 \(\mathrm{Ker}(u_{(\mathbf{C})}^* - i)\) 到 \(\mathrm{Ker}(u_{(\mathbf{C})}^* + i)\) 上的希尔伯特空间同构，于是可以应用推论 1。

推论 3. --- 设 u 是 E 上的对称算子。下列条件等价：

\begin{enumerate}
\def\labelenumi{(\roman{enumi})}
\item
  算子 u 本质自伴；
\item
  空间 \(\operatorname{Ker}(u^{*} + i)\) 与 \(\operatorname{Ker}(u^{*} - i)\) 为零；
\item
  空间 \(\operatorname{Im}(u + i)\) 与 \(\operatorname{Im}(u - i)\) 在 E 中稠密；
\item
  对称部分算子 u 有唯一的自伴扩张。
\end{enumerate}

由 IV，第 236 页命题 7，c)，条件 (ii) 与 (iii) 等价。

部分算子 \(u\) 本质自伴，当且仅当 \(\overline{u}\) 自伴，亦即当且仅当 \(\text{Ker}(\overline{u}^* - i) = \text{Ker}(\overline{u}^* + i) = \{0\}\)（命题 26 的推论）。由于 \(\overline{u}^* = u^*\)（IV，第 237 页，命题 8），条件 (i) 因此等价于条件 (ii)。此外，前一推论表明 (ii) 蕴含 \(u\) 有唯一的自伴扩张，此即条件 (iv)。

最后，设 \(u\) 有唯一的自伴扩张 \(v\)。那么 \(\overline{u}\) 亦然，因为 \(\overline{u}\) 的每个自伴扩张都是 \(u\) 的自伴扩张，从而等于 \(v\)。由前一推论，空间 \(\mathrm{Ker}(\overline{u}^{*} + i)\) 与 \(\mathrm{Ker}(\overline{u}^{*} - i)\) 必为零，再利用 \(\overline{u}^{*} = u^{*}\) 即得条件 (ii)。

例。--- 设 U 是 \(\mathbf{R}^{n}\) 中赋有勒贝格测度 \(\mu\) 的开子集。设 \(\Delta\) 是数量微分算子

\[
\Delta = - \sum_ {i = 1} ^ {n} \partial_ {i} ^ {2}
\]

此算子定义在 U 上。用 \(u\) 表示实希尔伯特空间 \(\mathrm{L}_{\mathbf{R}}^{2}(\mathrm{U},\mu)\) 上以 \(\mathcal{D}_{\mathbf{R}}(\mathrm{U})\)（IV，第 201 页，注）为定义域、由 \(\varphi \mapsto -\sum_{i=1}^{n}\partial_{i}^{2}\varphi\) 定义的部分算子。我们有 \(u_{(\mathbf{C})} = \Delta_{-}\)，它是可闭的部分算子（IV，第 242 页，命题 13）且对称（因为 \(\Delta\) 形式上对称），故 \(u\) 可闭且对称（IV，第 228 页命题 2 与 IV，第 239 页引理 3）。于是由上述推论，存在 \(u_{(\mathbf{C})}\) 的自伴扩张。这就是 U 上的一个拉普拉斯算子（IV，第 243 页）。

